\section{基于视觉语言模型的图像融合方法研究}\label{ch:method}

\subsection{引言}\label{sec:ch3_intro}

本章作为全文研究的核心方法部分，主要围绕基于视觉语言模型的多曝光图像融合方法FILM展开系统阐述。该部分承接第二章中关于深度学习图像融合、视觉语言模型、跨注意力机制以及多曝光图像处理基础的理论分析，进一步给出FILM方法的整体框架、关键模块实现、训练策略与实验验证，并为第四章交互式图像融合系统的设计与实现提供算法基础和性能依据。本章首先对FILM方法的问题定义、处理流程和整体架构进行说明，明确文本语义引导图像融合的基本思路；接着详细介绍文本语义特征提取与预处理、视觉-文本跨注意力融合块、多阶段级联融合结构以及损失函数与训练策略；最后通过定量实验、定性分析和模块作用分析验证方法在多曝光图像融合任务中的适用性，为后续系统级应用奠定基础。

\subsection{方法整体设计}\label{sec:method_overview}

\textbf{（1）问题定义与处理流程。}

本章详细阐述本文复现并适配的基于视觉语言模型的图像融合方法FILM（Image Fusion via Vision-Language Model）\supercite{zhao2024film}。FILM方法的核心思想是：利用视觉语言模型相关流程得到的高层文本语义特征，通过跨注意力机制对底层视觉特征的融合过程进行语义引导，使得融合结果在保留视觉细节的同时更加符合人类的语义理解。

FILM方法的处理流程如下：首先将过曝光和欠曝光图像的亮度通道输入模型，得到图像级描述、区域级描述及分割掩码三类信息，并借助 ChatGPT 将其转化为文本特征；随后用预训练的 BLIP‑2 将文本语义映射到与图像相同的向量空间，通过跨注意力机制引导两幅曝光图像的特征融合，再经三个 Restormer 模块精炼，得到融合后的亮度通道；色度通道采用等权重平均，最后合并 Y、Cb、Cr 并转回 RGB，输出最终的彩色融合图像。

\textbf{（2）FILM方法的整体架构。}

FILM方法的整体架构主要由以下三个部分组成：

（1）\textbf{文本语义特征提取模块}：使用预训练视觉语言模型相关流程得到的768维文本特征向量。BLIP2通过Q-Former模块在冻结的视觉编码器和大语言模型之间建立信息桥梁，能够生成高度语义化的文本嵌入表示；在本项目实现中，训练与测试阶段主要读取数据集中预先保存的文本特征。

（2）\textbf{视觉-文本跨注意力融合模块}：复现了一种基于Restormer Transformer的跨注意力块（restormer\_cablock），该模块对两幅源图像分别进行特征提取后，通过跨注意力机制将文本语义特征与视觉特征进行对齐和融合，实现文本语义对视觉特征处理的引导。

（3）\textbf{多阶段级联融合模块}：采用三级级联的跨注意力块结构，逐级深化文本语义引导下的特征融合。跨注意力块输出的特征经过拼接后，通过三个标准的Restormer Transformer块进行进一步精炼，最终通过1$\times$1卷积和Sigmoid激活函数生成单通道的融合亮度结果。

FILM方法的整体架构如\cref{fig:film_arch}所示。

\begin{figure}[htb!]
    \centering
    \includegraphics[width=0.95\textwidth]{img/fig3_1_film_arch.png}
    \caption{FILM方法整体架构}
    \label{fig:film_arch}
\end{figure}

\subsection{文本语义特征提取与预处理}\label{sec:text_preprocess}

\textbf{（1）BLIP2文本特征提取。}

预提取文本特征的维度为768，这些特征向量包含了源图像对应文本描述的高层语义信息。由于BLIP2使用冻结的大语言模型（如OPT或FlanT5），提取的文本特征具有稳定的语义表示能力，不需要在融合网络训练过程中进行微调。

\textbf{（2）文本特征投影与维度变换。}

提取的文本特征需要投影到与网络内部特征兼容的嵌入空间（hidden\_dim=256）。本文设计了一个简单的文本预处理模块（text\_preprocess），通过一维卷积实现维度变换：
\begin{equation}
    T' = \text{text\_process}(T) = \text{Conv1d}_{768 \rightarrow 256}(T^T)^T
\end{equation}
该模块使用一维卷积层，输入通道为768，输出通道为256，卷积核大小为1，步长为1，无填充。由于卷积核大小为1，该操作等价于对每个文本特征向量执行线性投影：
\begin{equation}
    T' = T W_t, \quad W_t \in \mathbb{R}^{768 \times 256}
\end{equation}
经过预处理后，文本特征$T' \in \mathbb{R}^{N \times 256}$可以在后续模块中与视觉特征进行跨注意力交互。

\subsection{视觉-文本跨注意力融合块}\label{sec:cross_attn_block}

\textbf{（1）Restormer特征提取。}

跨注意力融合块是FILM方法的核心模块，实现了文本语义引导下的视觉特征融合。该块对两幅源图像分别进行独立的特征处理。

对于源图像A，首先将输入图像$I_A$通过3$\times$3卷积和PReLU激活函数后，输入一个Restormer Transformer块进行特征提取：
\begin{equation}
    F_A = \text{Restormer}(\text{PReLU}(\text{Conv}_{3\times3}(I_A)))
\end{equation}
其中Restormer块使用多Dconv头转置自注意力（MDTA）和门控前馈网络（GDFN），能够同时捕捉全局和局部的视觉特征。

\textbf{（2）图像到文本特征投影。}

将提取的视觉特征投影到文本特征空间。该模块包含三个操作：首先通过1$\times$1卷积将特征通道数调整到image2text\_dim（32维）；然后通过最近邻插值将特征图的空间分辨率调整为288$\times$384；最后将空间特征展平为序列，得到与文本特征对齐的表示：
\begin{equation}
    F_{A2T} = \text{Reshape}(\text{Interpolate}(\text{Conv}_{1\times1}(F_A), [288, 384]))
\end{equation}
该投影操作使得视觉特征与文本特征在维度上完全对齐，为后续的跨注意力计算提供了前提条件。

\textbf{（3）跨注意力计算与语义引导。}

以预处理后的文本特征$T'$作为查询（Query），以投影后的视觉特征$F_{A2T}$作为键（Key）和值（Value），执行多头跨注意力计算：
\begin{equation}
    C_A = \text{CrossAttention}(T', F_{A2T}, F_{A2T}) = \text{softmax}\left(\frac{T' W_Q (F_{A2T} W_K)^T}{\sqrt{d_k}}\right)(F_{A2T} W_V)
\end{equation}
其中$W_Q, W_K, W_V$分别为查询、键、值的线性投影矩阵，$d_k = 256 / 8 = 32$为每个注意力头的维度。

跨注意力机制通过以文本特征作为查询、图像特征作为键和值，使得文本语义信息能够有选择地从视觉特征中提取与语义描述相符的信息，过滤掉语义不相关的噪声和冗余。

\textbf{（4）注意力池化与特征重构。}

对跨注意力输出序列在序列维度上执行自适应平均池化，将序列长度压缩为1，然后进行L1归一化：
\begin{equation}
    \hat{C}_A = \text{Normalize}_{L1}(\text{AdaptiveAvgPool1d}(C_A^T)^T)
\end{equation}
L1归一化确保了注意力权重的和为1，使得后续的加权操作具有明确的概率解释。

将归一化后的注意力权重与视觉特征相乘，然后投影回空间维度，将重构后的特征与原始特征进行拼接，通过1$\times$1卷积进行融合。源图像B的处理流程与A完全对称，共享相同的网络结构但使用独立的参数。

跨注意力融合块的设计具有以下优势：

（1）\textbf{语义引导的特征选择}：文本特征作为查询的跨注意力机制实现有选择的信息提取。

（2）\textbf{双分支独立处理}：对两幅源图像分别进行独立的特征提取和跨注意力融合，避免了早期融合导致的信息损失。

（3）\textbf{残差连接}：在特征重构阶段引入残差连接，确保原始视觉信息的有效传递。

\subsection{多阶段级联融合结构}\label{sec:cascade_fusion}

\textbf{（1）三级级联架构设计。}

FILM方法采用了三级级联的跨注意力块结构，逐级深化文本语义引导下的特征融合。

第一级跨注意力块（restormer\_ca1）接收原始的1通道亮度图像作为输入，提取初始的视觉特征并进行文本语义引导的融合。

第二级跨注意力块（restormer\_ca2）接收第一级的输出特征（mid\_channel=32维）作为输入，在更深层次上进行文本语义引导。

第三级跨注意力块（restormer\_ca3）接收第二级的输出特征作为输入，进一步细化融合特征表示。

三级级联结构使得网络能够在不同抽象层次上进行文本语义引导的特征融合，从底层的边缘、纹理特征到高层的语义、结构特征，逐步建立起完整的融合表征。

\textbf{（2）特征拼接与Restormer精炼。}

经过三级跨注意力块处理后，来自源图像A和B的特征被拼接在一起：
\begin{equation}
    F_{cat} = \text{Concat}(F_A^{(3)}, F_B^{(3)}, \text{dim}=1) \in \mathbb{R}^{B \times 64 \times H \times W}
\end{equation}
拼接后的特征通过三个标准的Restormer Transformer块进行进一步精炼。第一个Restormer块的输入通道为64（32+32），其余为32。

\textbf{（3）融合亮度图生成。}

最后的1$\times$1卷积将特征图压缩为单通道输出，通过Sigmoid激活函数生成[0, 1]范围内的融合亮度结果：
\begin{equation}
    F_{fused} = \sigma(\text{Conv}_{1\times1}(F_4))
\end{equation}
其中$\sigma$为Sigmoid激活函数。该结果作为融合后的Y通道参与后续的色彩空间重建。

\subsection{损失函数与训练策略}\label{sec:loss_training}

\subsubsection{梯度损失}\label{sec:gradient_loss}

梯度损失鼓励融合图像在边缘和纹理细节方面尽可能接近或超越源图像，确保融合结果保留两幅源图像的全部细节信息。该损失基于Sobel算子计算图像的梯度信息：
\begin{equation}
    \mathcal{L}_{gradient} = \frac{1}{MN}\sum_{i,j} \|G_{fused}(i,j) - \max(G_A(i,j), G_B(i,j))\|_1
\end{equation}
其中$G_A$、$G_B$和$G_{fused}$分别表示源图像A、源图像B和融合图像的梯度图。梯度图通过Sobel算子计算得到：
\begin{equation}
    G = \sqrt{(\text{Sobel}_x * I)^2 + (\text{Sobel}_y * I)^2}
\end{equation}
其中$\text{Sobel}_x$和$\text{Sobel}_y$分别为水平和垂直方向的Sobel卷积核，$*$表示卷积操作。

梯度损失采用源图像梯度的最大值作为参考，鼓励融合图像在每个位置都保留梯度较大的边缘信息。这种设计基于以下假设：对于多曝光图像融合，每个像素位置至少在一幅源图像中具有清晰的边缘信息，融合结果应该在所有位置上都保留清晰的边缘。

\subsubsection{强度损失}\label{sec:intensity_loss}

强度损失鼓励融合图像的像素强度尽可能接近源图像强度的最大值，确保融合结果具有足够的亮度和对比度：
\begin{equation}
    \mathcal{L}_{intensity} = \frac{1}{MN}\sum_{i,j} \|I_{fused}(i,j) - \max(I_A(i,j), I_B(i,j))\|_1
\end{equation}
其中$I_A$、$I_B$和$I_{fused}$分别表示源图像A、源图像B和融合图像的像素强度值。

强度损失采用L1范数计算，对异常值具有较强的鲁棒性。与L2范数相比，L1范数不会因为个别像素的巨大差异而产生过大的梯度，从而有利于训练的稳定性。强度损失的参考值采用源图像像素强度的最大值，确保融合结果在亮度方面能够充分保留两幅源图像的优势。

融合损失的总体表达式为：
\begin{equation}
    \mathcal{L}_{fusion} = \lambda_1 \mathcal{L}_{gradient} + \lambda_2 \mathcal{L}_{intensity}
\end{equation}
其中$\lambda_1 = 20$为梯度损失的权重系数，$\lambda_2 = 1$为强度损失的权重系数。梯度损失的权重远大于强度损失，体现了本文对边缘和纹理细节保留的重视程度高于单纯的亮度重建。

此外，系统还支持结构相似性（SSIM）损失作为可选的补充损失。SSIM损失通过衡量融合图像与源图像之间的结构相似性来指导训练：
\begin{equation}
    \mathcal{L}_{ssim} = 1 - \frac{\text{SSIM}(I_{fused}, I_A) + \text{SSIM}(I_{fused}, I_B)}{2}
\end{equation}
在本文的实验中，SSIM损失的权重系数默认设置为0，即不启用SSIM损失。这是因为SSIM损失在实验中发现对融合质量的提升有限，且会增加训练的计算开销。

\subsubsection{训练配置与策略}\label{sec:training_config}

FILM模型训练脚本支持使用PyTorch的DataParallel策略进行数据并行训练。主要的训练配置如下：

（1）\textbf{数据集}：使用VLFDataset数据集及其预处理后的HDF5文件。该数据集包含多曝光、红外-可见光、医学、多聚焦等多种融合任务的图像对和对应的文本特征。每个图像对包含两幅源图像（imageA和imageB）以及对应的文本特征向量。

（2）\textbf{优化器}：采用Adam优化器，初始学习率设为1$\times$10$^{-5}$。Adam优化器结合了动量和自适应学习率的优点，在深度学习训练中表现出良好的收敛性能。

（3）\textbf{批次大小}：训练脚本中DataLoader的批次大小默认设置为2。较小的批次大小有助于降低显存占用，同时便于在普通实验环境中完成模型训练或微调。

（4）\textbf{训练轮数}：总训练轮数为150个epoch。经过实验验证，150个epoch足以使模型收敛，同时避免过拟合。

（5）\textbf{学习率调度}：采用StepLR阶梯式学习率调度策略。学习率从初始值1$\times$10$^{-5}$开始，每隔\texttt{step\_size=50}个epoch乘以衰减系数\texttt{gamma=0.6}，使模型在训练后期以较小学习率继续优化。

（6）\textbf{数据预处理}：训练前通过data\_process.py脚本将原始图像和文本特征预处理为HDF5格式，以提高数据加载效率。HDF5文件包含imageA、imageB和text三个数据集组，分别存储两幅源图像和对应的文本特征。

（7）\textbf{检查点保存}：每个epoch结束后保存完整的模型检查点，包括网络状态字典、优化器状态和学习率调度器状态，以便训练中断后的恢复。

训练过程中，若实验环境包含多块GPU，模型权重可通过DataParallel机制在GPU之间共享并进行梯度聚合；若仅使用单块GPU，训练脚本同样可以正常执行。该设计提高了训练脚本在不同硬件环境下的适配性。

\subsection{实验与分析}\label{sec:experiments}

\subsubsection{实验设置}\label{sec:experiment_setup}

\textbf{数据集}：实验在MEFB（Multi-Exposure Fusion Benchmark）数据集上进行验证。该数据集包含了多组不同场景的多曝光图像对，涵盖了室内、室外、自然、建筑等多种类型。每组图像对包含一幅过曝光图像和一幅欠曝光图像，两幅图像拍摄的是同一场景。

\textbf{评价指标}：采用\cref{sec:metrics}中介绍的六种客观评价指标对融合结果进行全面评估，包括信息熵（EN）、标准差（SD）、空间频率（SF）、平均梯度（AG）、视觉信息保真度（VIF）和边缘信息保持度（Qabf）。

\textbf{对比方法}：实验选取了以下传统融合方法与前沿深度学习方法作为对比基准：

（1）\textbf{均值融合（Average）}：将两幅源图像的对应像素值进行等权重平均。该方法计算简单、速度快，但无法根据图像内容进行自适应的融合决策，融合结果往往对比度不足。

（2）\textbf{像素最大值融合（Max）}：取两幅源图像对应像素值的最大值作为融合结果。该方法能够在一定程度上保留亮度较高的区域细节，但在暗部区域容易出现信息丢失。

（3）\textbf{Mertens融合}\supercite{mertens2009exposure}：基于对比度、饱和度和曝光度三个质量度量的拉普拉斯金字塔融合方法。该方法在传统融合方法中表现优异，被广泛应用于消费级相机的HDR模式。本实验中通过OpenCV的\texttt{createMergeMertens}实现，对源图像的YCbCr亮度通道进行融合后提取Y通道进行指标评估。

（4）\textbf{FFMEF}\supercite{zheng2023ffmef}：CVPR Workshop 2023提出的高效多曝光图像融合方法。该方法采用滤波器主导的融合策略和梯度驱动的无监督学习，通过多尺度特征提取与空间交叉注意力机制实现自适应融合。与FILM不同，FFMEF仅依赖低层视觉特征（亮度、梯度），不引入文本语义引导。本实验使用其官方提供的MEF任务预训练权重进行推理。

\textbf{模型检查点}：使用预训练的MEF.pth模型权重进行推理，该模型面向多曝光图像融合任务训练。加载时需要移除DataParallel保存的"module."前缀，将模型设置到GPU或CPU上并切换为评估模式。

\textbf{推理流程}：融合推理过程中，首先将RGB源图像转换到YCbCr色彩空间，提取Y通道输入到FILM网络进行融合推理；对网络输出进行归一化处理，并将其作为融合后的Y通道；Cb和Cr通道采用等权重的经典加权平均进行融合；最后将融合后的Y、Cb、Cr通道合并，转换回RGB色彩空间输出结果。推理过程中，模型仅处理图像的Y（亮度）通道，既保证了融合质量又有效降低了计算成本。

\subsubsection{定量实验结果分析}\label{sec:quantitative_results}

FILM方法与对比方法的定量比较结果如\cref{tab:quantitative}所示。

\begin{table}[htb!]
    \caption{不同融合方法的客观评价指标对比（MEFB数据集，40张测试图像均值）}
    \label{tab:quantitative}
    \centering
    \begin{tabular}{ccccccc}
        \hline
        方法 & EN & SD & SF & AG & VIF & Qabf \\
        \hline
        Average & 6.80 & 50.35 & 11.16 & 3.23 & 1.44 & 0.59 \\
        Max & 5.90 & 64.32 & 16.49 & 4.30 & 1.45 & 0.69 \\
        Mertens & 6.80 & 51.41 & 13.64 & 3.61 & 1.12 & 0.47 \\
        FFMEF & 6.89 & 56.85 & 12.93 & 3.72 & 1.46 & 0.62 \\
        FILM & \textbf{7.32} & \textbf{68.97} & \textbf{20.96} & \textbf{6.11} & \textbf{1.52} & \textbf{0.77} \\
        \hline
    \end{tabular}
\end{table}

从\cref{tab:quantitative}的定量比较结果可以看出，FILM方法在全部六项评价指标上均优于传统融合方法和FFMEF深度学习方法。

在信息熵（EN）方面，FILM方法达到了7.32，高于所有对比方法（FFMEF为6.89，Average和Mertens均为6.80，Max为5.90）。信息熵衡量融合图像包含的信息量，FILM方法在保留源图像细节的同时保持了最丰富的灰度层次分布。

在标准差（SD）方面，FILM方法达到了68.97，说明融合图像具有较高的对比度和丰富的灰度层次分布。与FFMEF方法（56.85）相比，FILM方法的SD提高了约21.3\%，与Max方法（64.32）相比提高了约7.2\%，与Mertens方法（51.41）相比提高了约34.2\%，与Average方法（50.35）相比提高了约36.9\%。这一结果表明，文本语义引导使得网络能够更准确地保留源图像中的有效亮度与结构信息。

在空间频率（SF）和平均梯度（AG）方面，FILM方法分别取得了20.96和6.11的优异成绩，显著高于所有对比方法（FFMEF方法分别为12.93和3.72，Max方法分别为16.49和4.30，Mertens方法分别为13.64和3.61，Average方法分别为11.16和3.23）。特别是SF指标，FILM方法比FFMEF提高了约62.1\%，比Max方法提高了约27.1\%。这表明融合图像的空间活跃度和细节清晰度得到了有效提升，文本语义引导的网络成功地将两幅源图像的边缘和纹理信息融合到了最终结果中。

在视觉信息保真度（VIF）方面，FILM方法达到了1.52，高于FFMEF方法的1.46、Max方法的1.45和Average方法的1.44，也显著优于Mertens方法的1.12。VIF基于自然场景统计模型和小波分析，能够更准确地反映人眼对融合质量的主观感知。FILM方法在VIF指标上的表现说明其融合结果能够较好地保留源图像中的视觉信息。

在边缘信息保持度（Qabf）方面，FILM方法达到了0.77，表明融合图像较好地保留了源图像的边缘结构信息。与FFMEF方法（0.62）相比提高了约24.2\%，与Max方法（0.69）相比提高了约11.6\%，与Average方法（0.59）相比提高了约30.5\%，与Mertens方法（0.47）相比提高了约63.8\%。

综合来看，FILM方法在全部六项指标上均取得最优结果。与仅依赖低层视觉特征的FFMEF方法相比，FILM方法在SD、SF、AG和Qabf四项核心指标上的优势尤为明显，这验证了引入文本语义引导机制在多曝光图像融合任务中的有效性——高层语义信息能够帮助网络做出更加准确的融合决策。

\subsubsection{定性实验结果分析}\label{sec:qualitative_results}

从主观视觉效果来看，在不同场景下FILM方法均表现出优异的融合质量：

（1）\textbf{室内场景}：在包含窗户等强光源的室内场景中，FILM方法的融合结果同时保留了室内外的细节信息，窗框边缘清晰，室内物品的纹理和色彩还原自然，整体曝光均衡。均值融合方法的结果则显得灰暗，对比度不足；Max方法在过曝光区域出现了明显的亮度溢出。

（2）\textbf{室外自然场景}：FILM方法的融合结果同时保留了天空的云层纹理和地面的植被、建筑细节，整体色彩自然，过渡平滑。Mertens方法虽然也能较好地保留细节，但在天空与地面的交界处存在一定的色彩失真。FFMEF方法在整体曝光均衡性上优于传统方法，但在纹理细节保留方面略逊于FILM。

（3）\textbf{建筑场景}：FILM方法的融合结果在建筑细节和天空层次方面均表现优异，建筑立面的纹理、门窗等细节清晰可见，天空的蓝色渐变自然。传统方法在建筑暗面区域的细节恢复不够充分。

综合定量和定性实验结果，FILM方法在多曝光图像融合任务上取得了优于传统方法的整体融合性能，验证了文本语义引导机制和跨注意力融合设计在该任务中的适用性。

\subsubsection{模块作用分析与实现验证}\label{sec:ablation}

由于重新训练不同结构变体需要较高的计算资源，本项目未对w/o text、w/o cross-attn等结构变体进行完整重训练。因此，本文不将相关变体结果作为独立完成的消融实验，而是结合模型结构和代码实现对关键模块的作用进行分析。

\textbf{（1）文本语义引导模块。}

FILM方法通过文本特征为图像融合过程提供高层语义信息。在代码实现中，文本特征首先经过\texttt{text\_preprocess}模块由768维投影到256维，使其能够与视觉特征进入同一跨注意力计算空间。该设计使模型在处理多曝光图像时不仅依赖灰度、纹理和梯度等低层视觉信息，也能够利用场景描述中包含的语义线索辅助特征选择。

\textbf{（2）跨注意力融合模块。}

跨注意力融合块以文本特征作为Query，以图像特征作为Key和Value，实现文本到视觉特征的引导。与简单拼接不同，跨注意力能够根据文本语义动态聚合与当前语义相关的视觉区域。代码中两幅源图像分别经过独立分支处理，能够避免过早融合导致的细节损失；随后再通过特征拼接和Restormer精炼块完成最终融合。

\textbf{（3）三级级联结构。}

模型连续使用三个\texttt{restormer\_cablock}，使文本语义引导能够在不同特征层次上逐步发挥作用。第一级面向输入亮度图像提取初始结构信息，第二级和第三级在更深层特征空间中继续强化语义相关区域和纹理细节。该级联结构提升了模型对复杂曝光差异的表达能力，但也带来了更高的参数量和推理开销。

\textbf{（4）实现验证。}

从实现角度看，模型能够稳定完成文本特征预处理、双分支视觉特征提取、跨注意力交互、Restormer精炼和单通道融合亮度结果输出。结合前述定量实验可知，完整模型在全部六项指标上均优于传统融合方法和FFMEF深度学习方法，说明该结构在多曝光图像融合任务中具有较好的工程可用性。

\subsection{本章小结}\label{sec:ch3_summary}

本章详细阐述了本文复现并适配的基于视觉语言模型的图像融合方法FILM。首先介绍了方法的整体架构和设计思路，然后从文本特征预处理、跨注意力融合块、多阶段级联融合结构、损失函数和训练策略等多个层面进行了技术分析。实验结果表明，FILM方法在MEFB数据集上的全部六项评价指标（EN、SD、SF、AG、VIF、Qabf）均优于传统融合方法和FFMEF深度学习方法。与仅依赖低层视觉特征的FFMEF相比，FILM在SD、SF、AG和Qabf四项核心指标上的优势尤为明显，验证了文本语义引导机制在多曝光图像融合中的有效性。模块作用分析进一步说明了文本语义引导和跨注意力机制在融合过程中的作用。
